%% kapa: exam=1 round=21 year=2010 period=2 form=A subject=appraisal_law subject_ko=부동산관계법규 seq=8 number=48
%% source: Q-Net 감정평가사 기출문제 (cst00309) artlSeq=5024181 file=extracted/1008/5024181/2010년 제21회 감정평가사 1차시험 기출문제/2교시.hwp
%% conversion: kordoc 4.24.51 -> tools/gen1.py (auto); review=proofread-gpt-sol; flags=none
\begin{question}{48}
국토의 계획 및 이용에 관한 법령상 개발행위허가의 제한에 관한 설명으로 옳지 않은 것은?
\begin{choices}
  \item 국토해양부장관이 개발행위허가를 제한하고자 할 때는 중앙도시계획위원회의 심의를 거쳐야 한다.
  \item 기반시설부담구역으로 지정된 지역의 개발행위허가의 제한은 1회에 한하여 2년 이내의 기간 동안 할 수 있다.
  \item 시․도지사가 개발행위허가를 제한하고자 할 때는 시․도도시계획위원회의 심의 전에 미리 제한하고자 하는 지역을 관할하는 시장 또는 군수의 의견을 들어야 한다.
  \item 개발행위로 인하여 주변의 환경․경관․미관․문화재 등이 크게 오염되거나 손상될 우려가 있는 지역에 대한 개발행위허가의 제한은 1회에 한하여 3년 이내의 기간 동안 할 수 있다.
  \item 시․도지사의 개발행위 허가의 제한에 대한 고시는 당해 지방자치단체의 공보에 게재하는 방법으로 한다.
\end{choices}
\end{question}
\answer{2}
